\documentclass[11pt,a4paper]{article}
\usepackage{fontspec}
\usepackage{tikz}
\usepackage{tikz-uml}


%%\setmainfont[BoldFont=Adobe Heiti Std]{Adobe Song Std}
\setmainfont{SimSun}
%%\setsansfont[BoldFont=Adobe Heiti Std]{AR PL KaitiM GB}
%%\setmonofont{Bitstream Vera Sans Mono}
\begin{document}
世界,你好!

      sudo mkdir /usr/share/fonts/consolas
      sudo cp 存放字体的目录/consolas.ttf  /usr/share/fonts/consolas
      sudo chmod 644 /usr/share/fonts/consolas/* 

      cd /usr/share/fonts/consolas
      sudo mkfontscale
      sudo mkfontdir 
      sudo fc-cache -fv 

指令全名：set-buffer-file-coding-system
指令作用：改变当前buffer的编码
调用方法：(以目标编码gbk为例)
C-x <RET> f gbk <RET>
该指令还可以用于改变当前buffer的换行习惯(编码参数用dos或unix，对应\r\n或\n)。

指令全名：universal-coding-system-argument
指令作用：指定紧随其后的命令(如C-x C-f或C-x C-w等)所采用的编码
调用方法：(以目标编码gbk为例)
C-x <RET> c gbk <RET>

指令全名：revert-buffer-with-coding-system
指令作用：用指定编码重读当前buffer(如果打开时用错了编码)
调用方法：(以目标编码gbk为例)
C-x <RET> r gbk <RET>

指令全名：recode-region
指令作用：将以错误编码解码的选区以指定编码重新解码
调用方法：(以目标编码gbk为例)
M-x recode-region <RET> gbk <RET> utf-8 <RET>

除了这些操作外，还有一种方法，那就是在文件开始的部分给出如下形式的指令，直接告诉编辑器在打开和保存时应采用的编码：
-*- coding: gbk -*-
接触过Python的同学是不是觉得很眼熟?

[补充] 查看当前编码选择的命令为：
M-x describe-coding-system <RET> 或 C-h C <RET>
也可在回车前输入具体的编码名称(如gbk)以查看详细说明。

\begin{tikzpicture}
\draw (0,0) --(1,2) -- (2,3) -- (1,0);
\draw (-1,3) --(1.5, 0.5);
\end{tikzpicture}

\begin{tikzpicture} \draw [yellow, line width=6] (0,0) -- (.5,0); \end{tikzpicture} you want

\begin{tikzpicture}

\umlclass[x=0, y=-3] {namespace::class1\_name类1}{}{}
\umlclass[x=0, y=-7, width=35ex] {namespace::class2\_name}{}{}
\umlclass[x=8, y=-5] {namespace::class3\_name}{}{}
\umlclass[x=0, y=-10] {namespace::class4\_name} {}{}
\umluniaggreg [arg1=*, mult1=1, arg2=a, mult2=1] {namespace::class1\_name类1} {namespace::class2\_name}
\umluniaggreg [arg2=b, mult2=1,pos=0.95] {namespace::class1\_name类1} {namespace::class3\_name}
\umlinherit {namespace::class4\_name}{namespace::class3\_name}
\umlinherit {namespace::class4\_name}{namespace::class2\_name}

\begin{umlpackage} [x=0, y=0] {My First Package 包1}
\end{umlpackage}

\begin{umlpackage} [x=7, y=-3] {My Second Package 包2}
\end{umlpackage}
%%\end{tikzpicture}

\begin{umlsystem} [x=5,y=5] {use case system 用户例图}
%%\umlactor [scale=0.5] {small user}
\umlusecase [name=c2] {Case 2}
\umlusecase [x=1, y=-2, name=c3] {Case 3 事例3}
\umlusecase [y=-4, name=c5] {Case 5}
\umlinclude {c2} {c3}
\umlinherit {c5} {c3}
\end{umlsystem}
\umlactor [x=0, y=5, scale=0.7] {user}
\umlassoc {user}{c2}
\umlassoc {user} {c5}
\umlusecase [x=5, y=5, name=c1] { case 1}
%%\begin{tikzpicture}
\end{tikzpicture}

\end{document}
